\documentclass[11pt,a4paper]{article}
\usepackage[utf8]{inputenc}
\usepackage[french]{babel}
\usepackage{amsmath,amssymb,amsthm}
\usepackage{geometry}
\geometry{hmargin=2.5cm,vmargin=3cm}
\usepackage{tikz}
\usetikzlibrary{cd,positioning,shapes}
\usepackage{listings}
\usepackage{xcolor}

\lstset{
    language=Python,
    basicstyle=\ttfamily\small,
    keywordstyle=\color{blue}\bfseries,
    stringstyle=\color{red},
    commentstyle=\color{gray}\itshape,
    numbers=left,
    numberstyle=\tiny\color{gray},
    breaklines=true,
    frame=single,
    showstringspaces=false,
    literate={é}{{\'e}}1 {è}{{\`e}}1 {à}{{\`a}}1 {ê}{{\^e}}1 {î}{{\^i}}1 {ç}{{\c{c}}}1 {ï}{{\"i}}1 {É}{{\'E}}1,
    extendedchars=true,
    inputencoding=utf8
}

\title{\Huge DÉMONSTRATION DU THÉORÈME DE CONVERGENCE MONOTONE \\ \large Cours magistral, exercices corrigés et travaux pratiques}
\author{\Large Charles EDOU NZE}
\date{\today}

\begin{document}
\maketitle
\tableofcontents
\newpage

\part{Cours Magistral}


\section{Théorème de convergence monotone (Beppo Levi)}

\subsection{Présentation du concept clé}



\begin{figure}[h]
\centering
\begin{tikzpicture}
  \draw[->] (-1,0) -- (6,0) node[right] {$x$};
  \draw[->] (0,-0.5) -- (0,5) node[above] {$y$};

  \draw[domain=0:5, smooth, variable=\x, blue, thick] plot ({\x}, {4 - 4*exp(-\x)});
  \draw[domain=0:5, smooth, variable=\x, blue, dashed] plot ({\x}, {4});

  \foreach \i in {1, 2, 3, 4} {
    \draw[domain=0:5, smooth, variable=\x, red, thin] plot ({\x}, {(4 - 4*exp(-\x))*(1 - exp(-\i))});
  }

  \node[blue, right] at (5,4) {$f$};
  \node[red, right] at (5,3) {$f_n$};
\end{tikzpicture}
\caption{Convergence monotone : une suite croissante de fonctions vers une limite.}
\end{figure}
\subsection{Formalisation}

Soit $(X, \mathcal{F}, \mu)$ un espace mesuré.

\subsubsection{Énoncé du Théorème}

\textbf{Théorème de Convergence Monotone (Beppo Levi) :}
Soit $(f_n)_{n \in \mathbb{N}}$ une suite de fonctions mesurables de $X$ dans $[0, +\infty]$.
Si la suite est \textbf{croissante} presque partout :
$$\forall n \in \mathbb{N}, \quad f_n \le f_{n+1} \text{ p.p.}$$
Alors la fonction limite $f = \lim_{n \to \infty} f_n$ est mesurable et :
$$\int_X \left( \lim_{n \to \infty} f_n \right) d\mu = \lim_{n \to \infty} \int_X f_n d\mu$$

\subsubsection{Corollaire (Sommation terme à terme)}

\textbf{Théorème :} Pour toute suite de fonctions mesurables \textbf{positives} $(u_n)$ :
$$\int_X \left( \sum_{n=0}^\infty u_n \right) d\mu = \sum_{n=0}^\infty \int_X u_n d\mu$$


\textbf{Exemple concret 1 : Suite géométrique}
Soit $u_n(x) = x^n$ sur $X = [0,1]$. Les $u_n$ sont positives et mesurables.
Le théorème permet d'écrire $\int_0^1 \sum_{n=0}^\infty x^n dx = \sum_{n=0}^\infty \int_0^1 x^n dx$.
$\int_0^1 x^n dx = \frac{1}{n+1}$. La somme de ces intégrales est la série harmonique $\sum \frac{1}{n+1} = +\infty$.
La somme de la série de fonctions est $\frac{1}{1-x}$, dont l'intégrale sur $[0,1]$ vaut aussi $+\infty$.

\textbf{Exemple concret 2 : Série à convergence très rapide}
Soit $u_n(x) = \frac{x^n}{n!}$ sur $X = [0,1]$.
$\int_0^1 \sum_{n=0}^\infty \frac{x^n}{n!} dx = \int_0^1 e^x dx = e - 1$.
Par le théorème, ceci est égal à $\sum_{n=0}^\infty \int_0^1 \frac{x^n}{n!} dx = \sum_{n=0}^\infty \frac{1}{n!(n+1)} = \sum_{n=0}^\infty \frac{1}{(n+1)!} = e - 1$.



\subsection{Démonstrations}

\subsubsection{Démonstration du Théorème de Beppo Levi}

1. \textbf{Existence de la limite :} Comme $(f_n(x))$ est une suite croissante de $[0, +\infty]$, elle admet toujours une limite dans $[0, +\infty]$ pour chaque $x$. On a vu (Jalon 65) que le sup (ou la limite ici) de fonctions mesurables est mesurable.
2. \textbf{Inégalité facile ($\ge$) :} Comme $f_n \le f$ pour tout $n$, par croissance de l'intégrale (Jalon 66) :
   $\int f_n \le \int f$. En prenant la limite : $\lim \int f_n \le \int f$.
3. \textbf{Inégalité difficile ($\le$) :} Soit $s$ une fonction simple telle que $0 \le s \le f$. Soit $\alpha \in ]0, 1[$.
   On définit $A_n = \{x \in X \mid f_n(x) \ge \alpha s(x) \}$.
   - Comme $f_n$ croît vers $f$ et $\alpha s < f$ (là où $s>0$), la suite d'ensembles $(A_n)$ est croissante et son union est $X$.
   - On a $\int f_n \ge \int_{A_n} f_n \ge \int_{A_n} \alpha s = \alpha \int_{A_n} s$.
   - Par continuité monotone de la mesure (Jalon 63), $\lim \int_{A_n} s = \int_X s$.
   - Donc $\lim \int f_n \ge \alpha \int_X s$.
   - En faisant tendre $\alpha \to 1$, on a $\lim \int f_n \ge \int s$.
4. \textbf{Conclusion :} Comme c'est vrai pour tout $s \le f$, alors $\lim \int f_n \ge \sup \int s = \int f$.
   Les deux inégalités prouvent l'égalité.

\subsection{Exercices d'Application}

\subsubsection{Exercice 1 : Intégrale d'une série}
\textbf{Énoncé :} Calculer $\int_0^1 \sum_{n=1}^\infty x^n dx$.
\textbf{Correction Détaillée :}
1. Les fonctions $u_n(x) = x^n$ sont mesurables et positives sur $[0, 1]$.
2. Par le corollaire du TCM, on peut intervertir somme et intégrale.
3. $I = \sum_{n=1}^\infty \int_0^1 x^n dx = \sum_{n=1}^\infty \left[ \frac{x^{n+1}}{n+1} \right]_0^1 = \sum_{n=1}^\infty \frac{1}{n+1}$.
4. Cette série est la série harmonique (moins son premier terme), elle diverge vers $+\infty$.
5. L'intégrale de la somme est donc $+\infty$.

\subsubsection{Exercice 2 : Niveau Avancé (Utilisation de la mesure de comptage)}
\textbf{Énoncé :} Retrouver le théorème de sommation des séries à termes positifs doubles : $\sum_{i} \sum_{j} a_{i,j} = \sum_{j} \sum_{i} a_{i,j}$.
\textbf{Correction Détaillée :}
On considère l'espace $\mathbb{N}$ muni de la mesure de comptage $\mu$. Soit $f_n(i) = \sum_{j=0}^n a_{i,j}$. La suite de fonctions $(f_n)$ est croissante car $a_{i,j} \ge 0$. Par Beppo Levi, l'intégrale de la limite est la limite des intégrales, ce qui correspond exactement à l'interversion des sommes.

\subsection{Application en Intelligence Artificielle}

-  En IA, on manipule souvent des \textbf{Espérances de sommes infinies} ou des limites de processus. Le TCM est l'outil qui permet de passer à la limite sous l'espérance en toute sécurité.
-
    - \textbf{Processus de Poisson :} Pour calculer le nombre moyen d'événements (ex: clics sur une pub) sur un intervalle de temps, on somme les probabilités d'événements infinitésimaux. Le TCM garantit que la somme de ces moyennes locales donne bien la moyenne globale.
    - \textbf{Séries de Taylor de fonctions de perte :} Si on décompose une fonction de coût complexe en une série de fonctions positives, on peut intégrer cette série terme à terme pour obtenir une approximation de la perte attendue.
    - \textbf{Théorie des Noyaux (Kernels) :} De nombreux noyaux (comme le noyau RBF) peuvent être vus comme des sommes infinies de caractéristiques. Le TCM permet de manipuler ces représentations de dimension infinie comme si elles étaient finies lors des calculs d'intégrales de risque.

\subsection{Liens Sémantiques}

- \textbf{Concepts Précédents requis :} [[Jalon 66 (Construction de l'intégrale de Lebesgue pour les fonctions mesurables positives.).md]], [[Jalon 63 (Définition axiomatique d'une mesure).md]]
- \textbf{Concepts Futurs dépendants :} [[Jalon 68 (Lemme de Fatou et définition de l'intégrale pour les fonctions de signe quelconque).md]], [[Jalon 69 (Démonstration complète du théorème de convergence dominée de Lebesgue.).md]]

\part{Exercices d'Application et de Concours}
\subsection*{Exercice 1 : Étude d'une intégrale par limite croissante $\bigstar$$\bigstar$$\star$$\star$$\star$}

\textbf{Énoncé :}
Soit $f_n(x) = \left(1 - \frac{x}{n}\right)^n \chi_{[0, n]}(x)$ sur $\mathbb{R}^+$. Montrer que $f_n$ est une suite croissante de fonctions mesurables. En déduire $\lim_{n \to \infty} \int_0^n \left(1 - \frac{x}{n}\right)^n dx$.

\textbf{Correction :}
1. On a $f_n(x) = \exp(n \ln(1 - x/n)) \chi_{[0, n]}(x)$. Posons $g_n(x) = n \ln(1 - x/n)$. On calcule la dérivée par rapport à $n$ (ou on étudie le rapport) et on obtient que $f_n(x)$ est strictement croissante en $n$.
2. La limite ponctuelle est $f(x) = \lim \exp(n \ln(1 - x/n))$. Or $n \ln(1 - x/n) = n (-x/n + O(1/n^2)) = -x + O(1/n)$. Donc $f(x) = e^{-x}$.
3. Les $f_n$ sont positives et mesurables (continues sur leur support). Par le théorème de convergence monotone, $\lim \int_0^\infty f_n = \int_0^\infty \lim f_n = \int_0^\infty e^{-x} dx$.
4. L'intégrale de $e^{-x}$ sur $\mathbb{R}^+$ vaut $[-e^{-x}]_0^\infty = 1$. Donc la limite de l'intégrale est $1$.


\subsection*{Exercice 2 : Série entière et intégration $\bigstar$$\bigstar$$\bigstar$$\star$$\star$}

\textbf{Énoncé :}
On considère l'intégrale $I = \int_0^1 \frac{\ln(x)}{x-1} dx$. Montrer que $I = \sum_{n=1}^\infty \frac{1}{n^2}$ en utilisant un développement en série et le théorème de convergence monotone.

\textbf{Correction :}
1. Sur $]0, 1[$, on a $\frac{1}{1-x} = \sum_{n=0}^\infty x^n$.
2. Donc $\frac{\ln(x)}{x-1} = \frac{-\ln(x)}{1-x} = \sum_{n=0}^\infty (-x^n \ln(x))$.
3. Les fonctions $u_n(x) = -x^n \ln(x)$ sont strictement positives sur $]0, 1[$ et mesurables.
4. Par le corollaire du théorème de convergence monotone pour les séries (Beppo Levi), on peut intervertir la somme et l'intégrale : $\int_0^1 \sum u_n(x) dx = \sum \int_0^1 u_n(x) dx$.
5. Calculons $\int_0^1 -x^n \ln(x) dx$. Par intégration par parties : on pose $u = \ln(x), dv = -x^n dx$, d'où $du = dx/x, v = -x^{n+1}/(n+1)$.
L'intégrale vaut $[-\ln(x) x^{n+1}/(n+1)]_0^1 - \int_0^1 \frac{-x^n}{n+1} dx = 0 + [\frac{x^{n+1}}{(n+1)^2}]_0^1 = \frac{1}{(n+1)^2}$.
6. On obtient donc $\int_0^1 \frac{\ln(x)}{x-1} dx = \sum_{n=0}^\infty \frac{1}{(n+1)^2} = \sum_{k=1}^\infty \frac{1}{k^2} = \frac{\pi^2}{6}$.


\subsection*{Exercice 3 : Application aux densités de probabilité $\bigstar$$\bigstar$$\bigstar$$\bigstar$$\star$}

\textbf{Énoncé :}
Soit $X$ une variable aléatoire positive admettant une densité $f$. Montrer par le théorème de convergence monotone que $\mathbb{E}[X] = \int_0^\infty \mathbb{P}(X > t) dt$.

\textbf{Correction :}
1. Par définition, $\mathbb{E}[X] = \int_0^\infty x f(x) dx$. On peut écrire $x = \int_0^x 1 dt$.
2. Ainsi $\mathbb{E}[X] = \int_0^\infty \left( \int_0^\infty \chi_{\{t < x\}}(t) dt \right) f(x) dx$.
3. La fonction $(t, x) \mapsto \chi_{\{t < x\}}(t) f(x)$ est positive et mesurable sur $\mathbb{R}^+ \times \mathbb{R}^+$.
4. On peut appliquer le théorème de Tonelli (qui repose fondamentalement sur Beppo Levi). On intervertit l'ordre d'intégration :
$\mathbb{E}[X] = \int_0^\infty \left( \int_0^\infty \chi_{\{t < x\}}(t) f(x) dx \right) dt$.
5. L'intégrale interne est $\int_t^\infty f(x) dx = \mathbb{P}(X > t)$.
6. Donc $\mathbb{E}[X] = \int_0^\infty \mathbb{P}(X > t) dt$.


\subsection*{Exercice 4 : Contre-exemple sans la croissance $\bigstar$$\bigstar$$\bigstar$$\bigstar$$\bigstar$}

\textbf{Énoncé :}
Considérons l'espace $(\mathbb{R}, \mathcal{B}(\mathbb{R}), \lambda)$. Soit $f_n(x) = \frac{1}{n} \chi_{[0, n]}(x)$. La suite $(f_n)$ est-elle croissante ? Calculer la limite de l'intégrale et l'intégrale de la limite. Conclure.

\textbf{Correction :}
1. $f_n(x)$ prend la valeur $1/n$ sur $[0, n]$ et $0$ ailleurs. $f_{n+1}(x)$ prend la valeur $1/(n+1)$ sur $[0, n+1]$. Sur $[0, n]$, $f_n(x) = 1/n > 1/(n+1) = f_{n+1}(x)$. La suite n'est donc \textbf{pas} croissante (elle est décroissante sur les compacts).
2. La limite ponctuelle est $f(x) = \lim_{n \to \infty} \frac{1}{n} \chi_{[0, n]}(x) = 0$. Donc $\int_{\mathbb{R}} f(x) dx = 0$.
3. Calculons l'intégrale de chaque terme : $\int_{\mathbb{R}} f_n(x) dx = \int_0^n \frac{1}{n} dx = 1$.
4. La limite des intégrales est $\lim 1 = 1$.
5. On a $\int \lim f_n = 0 \neq 1 = \lim \int f_n$. Le théorème de convergence monotone ne s'applique pas car l'hypothèse de croissance est indispensable pour éviter la 'fuite de masse' à l'infini.


\subsection*{Exercice 5 : Passage de Riemann à Lebesgue $\bigstar$$\star$$\star$$\star$$\star$}

\textbf{Énoncé :}
Prouver que la fonction de Dirichlet $f(x) = 1$ si $x \in \mathbb{Q} \cap [0,1]$, et $0$ sinon, n'est pas intégrable au sens de Riemann mais l'est au sens de Lebesgue via une limite de fonctions étagées.

\textbf{Correction :}
1. Au sens de Riemann, toute somme de Darboux supérieure vaut $1$ et toute somme inférieure vaut $0$ car $\mathbb{Q}$ et $\mathbb{R} \setminus \mathbb{Q}$ sont denses dans $[0,1]$. L'intégrale de Riemann n'existe pas.
2. Au sens de Lebesgue, on peut énumérer les rationnels de $[0,1]$ : $q_1, q_2, \dots$. Soit $f_n(x) = \chi_{\{q_1, \dots, q_n\}}(x)$.
3. Les $f_n$ sont des fonctions étagées positives, et la suite est croissante.
4. $\int f_n d\lambda = \sum_{k=1}^n \lambda(\{q_k\}) = \sum 0 = 0$.
5. La limite de $f_n$ est exactement $f$. Par Beppo Levi, $\int f d\lambda = \lim \int f_n d\lambda = 0$. La fonction est donc intégrable au sens de Lebesgue et son intégrale est nulle.


\subsection*{Exercice 6 : Inégalité de Fatou stricte $\bigstar$$\bigstar$$\star$$\star$$\star$}

\textbf{Énoncé :}
Donner un exemple explicite d'une suite de fonctions positives $f_n$ telle que $\int \lim\inf f_n < \lim\inf \int f_n$.

\textbf{Correction :}
1. Le lemme de Fatou énonce que $\int \lim\inf f_n \le \lim\inf \int f_n$. Cherchons un cas d'inégalité stricte.
2. On prend l'espace $\mathbb{R}$ avec la mesure de Lebesgue. Soit la 'bosse glissante' : $f_n(x) = \chi_{[n, n+1]}(x)$.
3. Pour tout $x \in \mathbb{R}$, pour $n$ assez grand ($n > x$), $f_n(x) = 0$. Donc $\lim_{n \to \infty} f_n(x) = 0$. La limite inférieure est la fonction nulle $0$.
4. Son intégrale est $\int 0 d\lambda = 0$.
5. D'autre part, pour tout $n$, $\int f_n d\lambda = \lambda([n, n+1]) = 1$. Donc $\lim\inf \int f_n = 1$.
6. On a bien $0 < 1$. L'inégalité stricte se produit à cause de la perte de masse vers l'infini (non compacité). Ce problème n'apparaît jamais avec des suites croissantes.


\subsection*{Exercice 7 : Application à la série harmonique $\bigstar$$\bigstar$$\bigstar$$\star$$\star$}

\textbf{Énoncé :}
Etudier $\int_0^1 \sum_{n=1}^\infty x^n dx$ en justifiant chaque étape.

\textbf{Correction :}
1. On a une série de fonctions $u_n(x) = x^n$.
2. Sur l'intervalle $[0, 1]$, pour tout $n \ge 1$, $u_n(x) \ge 0$. Les fonctions sont mesurables car continues.
3. On peut appliquer le théorème de Beppo Levi pour intervertir somme et intégrale.
4. $\int_0^1 \sum_{n=1}^\infty x^n dx = \sum_{n=1}^\infty \int_0^1 x^n dx$.
5. On calcule l'intégrale : $\int_0^1 x^n dx = \left[ \frac{x^{n+1}}{n+1} \right]_0^1 = \frac{1}{n+1}$.
6. Donc l'intégrale de la somme vaut $\sum_{n=1}^\infty \frac{1}{n+1} = \frac{1}{2} + \frac{1}{3} + \frac{1}{4} + \dots$.
7. Il s'agit de la série harmonique tronquée, qui diverge vers $+\infty$. L'intégrale vaut donc $+\infty$.


\subsection*{Exercice 8 : Continuité d'une intégrale paramétrée $\bigstar$$\bigstar$$\bigstar$$\bigstar$$\star$}

\textbf{Énoncé :}
Soit $f(t) = \int_0^\infty e^{-tx} \frac{\sin^2(x)}{x^2} dx$. Prouver la continuité de $f$ sur $\mathbb{R}^+$

\textbf{Correction :}
1. On remarque que $f_n(x)$ n'est pas croissante en $n$. Pour la continuité, le théorème de convergence monotone seul ne suffit pas en général, mais on peut exprimer la différence.
2. Cependant, pour l'étude de la continuité à droite en 0, si on prend $t_n \downarrow 0$, alors la suite de fonctions $g_n(x) = e^{-t_n x} \frac{\sin^2(x)}{x^2}$ est bien \textbf{croissante}.
3. En effet, $t_n$ décroît, donc $-t_n x$ croît pour $x \ge 0$, donc $e^{-t_n x}$ croît. La fonction $\frac{\sin^2(x)}{x^2}$ est positive.
4. Les $g_n$ sont positives, mesurables et croissent vers $g(x) = 1 \cdot \frac{\sin^2(x)}{x^2}$.
5. Par le TCM, $\lim \int g_n = \int \lim g_n$, donc $\lim_{n \to \infty} f(t_n) = f(0)$.
6. Cela démontre la continuité à droite en 0.


\subsection*{Exercice 9 : Lien avec la mesure produit $\bigstar$$\bigstar$$\bigstar$$\bigstar$$\bigstar$}

\textbf{Énoncé :}
Démontrer que pour des fonctions mesurables positives étagées, l'intégrale double vérifie Beppo Levi.

\textbf{Correction :}
1. Soit $(f_n(x,y))$ une suite croissante de fonctions étagées positives sur $X \times Y$.
2. On s'intéresse à $I_n = \int_{X \times Y} f_n(x,y) d(\mu \otimes \nu)$.
3. Par définition de l'intégrale des fonctions étagées, c'est une somme finie de mesures de rectangles.
4. La croissance de la suite assure que $\lim I_n$ existe dans $[0, +\infty]$.
5. La limite de $f_n$ est mesurable positive $f(x,y)$. Par construction de l'intégrale de Lebesgue, $f$ est la limite d'une suite de fonctions étagées $s_m \le f$. Le supremum de $\int s_m$ est $\int f$. La correspondance stricte est assurée par Beppo Levi sur l'espace produit.


\subsection*{Exercice 10 : Démonstration du théorème de Beppo Levi $\bigstar$$\star$$\star$$\star$$\star$}

\textbf{Énoncé :}
Restituer les grandes lignes de la preuve du théorème de Beppo Levi.

\textbf{Correction :}
1. L'inégalité $\lim \int f_n \le \int f$ est évidente car $f_n \le f$, donc par monotonie de l'intégrale $\int f_n \le \int f$, d'où le passage à la limite.
2. Pour l'autre inégalité, on prend une fonction étagée $0 \le s \le f$ et un réel $0 < \alpha < 1$.
3. On définit les ensembles mesurables $A_n = \{x \mid f_n(x) \ge \alpha s(x)\}$.
4. Comme $f_n$ croît vers $f$ et $\alpha < 1$, l'union croissante des $A_n$ couvre tout l'espace $X$.
5. $\int f_n \ge \int_{A_n} f_n \ge \alpha \int_{A_n} s$.
6. La continuité séquentielle croissante de la mesure implique que $\lim \int_{A_n} s = \int_X s$.
7. En faisant $\alpha \to 1$ puis le supremum sur toutes les $s \le f$, on obtient $\lim \int f_n \ge \int f$.



\part{Travaux Pratiques et Simulations Algorithmiques}
\subsection*{TP 1 : Vérification numérique de la convergence monotone}

\begin{lstlisting}
def verif_tcm(n_max):
    # On simule la suite f_n(x) = (1 - x/n)^n sur [0, n] et sa limite exp(-x)
    import math
    dx = 0.01

    # Integrale limite exp(-x) sur [0, 50]
    limit_integral = 0
    x = 0
    while x < 50:
        limit_integral += math.exp(-x) * dx
        x += dx

    print(f"Intégrale de la limite (théorique ~ 1) : {limit_integral}")

    # Calcul des intégrales de f_n
    for n in [5, 10, 50, 100]:
        integral_fn = 0
        x = 0
        while x < n:
            integral_fn += ((1 - x/n)**n) * dx
            x += dx
        print(f"Intégrale de f_{n} : {integral_fn}")

verif_tcm(100)
\end{lstlisting}


\subsection*{TP 2 : Simulation de la bosse glissante (Perte de masse)}

\begin{lstlisting}
def bosse_glissante():
    # Simulation de f_n(x) = 1 sur [n, n+1], 0 ailleurs
    # L'intégrale vaut toujours 1
    # Mais la limite simple est 0
    dx = 0.1
    for n in [10, 50, 100]:
        integral = 0
        # support est [n, n+1]
        x = n
        while x <= n + 1:
            integral += 1.0 * dx
            x += dx
        print(f"n={n}, Integrale={integral}")

    print("La limite des intégrales est 1. La limite des fonctions est 0, d'intégrale 0.")

bosse_glissante()
\end{lstlisting}


\subsection*{TP 3 : Interversion série et intégrale}

\begin{lstlisting}
def interversion_serie():
    # Intégrale de sum_{k=1}^n x^k sur [0,1]
    dx = 0.001

    for n in [5, 10, 20]:
        integral = 0
        x = 0
        while x < 1:
            val = sum([x**k for k in range(1, n+1)])
            integral += val * dx
            x += dx

        theoric_sum = sum([1/(k+1) for k in range(1, n+1)])
        print(f"n={n}, Integrale={integral:.4f}, Somme_theorique={theoric_sum:.4f}")

interversion_serie()
\end{lstlisting}


\subsection*{TP 4 : Suite de fonctions étagées}

\begin{lstlisting}
def etagees_approximation():
    # On approche f(x) = x^2 par des fonctions étagées constantes sur des intervalles de taille 1/n
    # sur l'intervalle [0, 1]
    import math

    def f(x): return x**2

    for n in [4, 10, 100]:
        integral = 0
        pas = 1.0 / n
        for k in range(n):
            inf = k * pas
            sup = (k+1) * pas
            val_min = f(inf) # min sur [inf, sup] car f croissante
            integral += val_min * pas

        print(f"n={n}, Integrale inferieure={integral:.5f}")

    print("Limite theorique = 1/3 = 0.33333")

etagees_approximation()
\end{lstlisting}


\subsection*{TP 5 : Concentration de Dirac}

\begin{lstlisting}
def suite_dirac():
    # f_n(x) = n/2 sur [-1/n, 1/n]
    # L'intégrale reste 1 (pas de croissance globale, donc pas TCM)
    dx = 0.0001

    for n in [10, 100, 1000]:
        integral = 0
        x = -1.0/n
        while x <= 1.0/n:
            integral += (n/2.0) * dx
            x += dx
        print(f"n={n}, Integrale={integral:.4f}")

    print("L'intégrale vaut 1, mais la limite de la fonction est 0 presque partout (sauf en 0 ou elle vaut l'infini).")

suite_dirac()
\end{lstlisting}



\end{document}
